\section{Introduction}

\emph{Six Birds Theory} (SBT) is a closure-based emergence calculus introduced in~\cite{Tsiokos2026SixBirds}; readers unfamiliar with its primitives will find all required definitions self-contained in Section~\ref{sec:sbt}. This paper presents original computational results, mechanized proofs, and reproducible experiments that instantiate the SBT packaging language in the quantum domain.

Quantum mechanics is our most accurate framework for microscopic phenomena, yet its standard presentation invites persistent conceptual puzzles: the state vector appears to collapse discontinuously upon measurement; entanglement correlations are often described as ``spooky action at a distance''; and Schr\"odinger's cat seems to demand macroscopic superpositions of mutually exclusive realities. Much of the interpretational debate concerns what the quantum state \emph{is}---a real physical field, an element of reality, or a bookkeeping device---but typically proceeds within a formalism that conflates causal claims about the world with inferential bookkeeping about what can be recorded.

Following Spekkens, we take seriously the possibility that this is a category mistake.\cite{Spekkens2007Toy} The quantum state is a tool for prediction under limited access; elevating it to an ontic object tends to force ``surplus'' structure (e.g., influences that cannot be used to signal and that never appear in any admissible record). This clashes with a Leibnizian methodological principle: if two scenarios are empirically indiscernible at a given descriptive layer, then they should not be treated as distinct objects of reality at that layer.

We argue that SBT provides a suitable replacement language.\cite{Tsiokos2026SixBirds} It explicitly separates causal substrate dynamics from \emph{packaging} operations that stabilize some distinctions into record-level objects. In this view, ``collapse'' is not a dynamical discontinuity but an idempotent packaging map (a closure) induced by a chosen record algebra; incompatibility of measurement contexts manifests as noncommuting closures (route mismatch), not superluminal causation; and the cat paradox is localized: ``alive vs.\ dead'' becomes an object-level distinction only when a stable record is packaged.\footnote{Reproducibility: all figures and quoted numbers are generated by \texttt{python -m sbtq.run\_all --seed 0 --out artifacts}; see Appendix~\ref{app:repro}.}

\paragraph{Contributions.}
This paper contributes:
\begin{itemize}
  \item A Spekkens-driven reformulation of ``layer realism'' via a Leibniz quotient: empirically indiscernible microdescriptions are identified as the same object at the layer (mechanized in Lean).
  \item A closure-based account of objecthood in which packaging maps are idempotent and objects correspond to fixed points; contextual incompatibility appears as noncommuting packaging routes (route mismatch), with certified finite witnesses (Lean).
  \item A quantum instantiation in which dephasing in a record basis plays the role of packaging: it is an idempotent closure whose fixed points are exactly diagonal (record-classical) states (Lean), and whose noncommutation across bases is quantitatively visible (Python).
  \item Fully reproducible experiments that make the mechanism concrete: double slit with environment overlap, quantum eraser via conditional packaging, quantified route mismatch, the cat (global purity vs.\ packaged mixture and idempotent packaging), staged objecthood in a metastable Markov chain, and no-signalling vs.\ conditioning in an EPR pair (Python).
\end{itemize}

\paragraph{Computational contributions.}
From a computer science and information theory perspective, this paper delivers:
\begin{itemize}
  \item Mechanized structural lemmas in Lean~4: Leibniz quotient factorization, saturation idempotence, dephasing fixed-point characterization, finite total-variation data-processing inequality, and certified noncommutation witnesses.
  \item A reproducible Python simulation suite (\texttt{sbtq}) with deterministic seeding, hash-traceable artifacts, and a robustness sweep over ten seeds.
  \item Information-theoretic diagnostics (route mismatch via trace distance, audit contraction, visibility curves) as operational proxies for context dependence and record stability.
  \item A Markov-kernel packaging analogue demonstrating timescale-dependent idempotence defect in a purely classical substrate.
\end{itemize}

\paragraph{Code availability.}
The repository is available at:
\begin{itemize}
  \item \url{https://github.com/ioannist/six-birds-quantum}
\end{itemize}

\paragraph{Paper map.}
Section~\ref{sec:spekkens} introduces the Spekkens diagnosis and the Leibnizian layer principle. Section~\ref{sec:sbt} presents the SBT packaging language and fixed-point view of objecthood. Section~\ref{sec:qm-package} maps quantum mechanics into this package and quantifies route mismatch. Section~\ref{sec:doubleslit} treats the double-slit and quantum eraser as objecthood budgeting. Section~\ref{sec:cat} analyzes measurement and the cat in an explicit system--apparatus--environment model. Section~\ref{sec:markov} gives a classical metastable Markov analogue. Section~\ref{sec:no-go} reframes selected no-go pressures in terms of globally compatible packaging assumptions. Section~\ref{sec:discussion} summarizes limitations and future directions.
